\documentclass[11pt,letterpaper]{article}

\usepackage[T1]{fontenc}
\usepackage{amsmath,amssymb}
\usepackage{tikz}
\usepackage{pgfplots}
\pgfplotsset{compat=1.18}
\usepackage{algorithm}
\usepackage{algpseudocode}
\usepackage{booktabs}
\usepackage{enumitem}
\usepackage[hidelinks]{hyperref}

\newcommand{\N}{\mathbb{N}}
\newcommand{\Z}{\mathbb{Z}}
\newcommand{\Q}{\mathbb{Q}}
\newcommand{\R}{\mathbb{R}}
\newcommand{\C}{\mathbb{C}}
\newcommand{\Map}[3]{#1 \colon #2 \to #3}
\newcommand{\Mod}[1]{#1\text{-}\mathrm{Mod}}
\newcommand{\GroupHom}{\mathrm{Hom}_{\mathbf{Grp}}}
\newcommand{\RingHom}{\mathrm{Hom}_{\mathbf{Ring}}}

\newtheorem{theorem}{Theorem}
\newtheorem{lemma}[theorem]{Lemma}
\newtheorem{corollary}[theorem]{Corollary}
\newtheorem{definition}[theorem]{Definition}
\newtheorem{remark}[theorem]{Remark}
\newenvironment{proof}{\noindent\textit{Proof.}}{\hfill\(\square\)}

\title{Carry energy is zero or greater than one half}
\author{}
\date{\today}

\begin{document}
\maketitle

\begin{abstract}
For odd positive integers \(q\) and \(c\), a finite base-two parity word whose odd letters are strictly contracting has a rational fixed point. The carry energy of the word is the greatest distance from a phase of that fixed point to the nearest positive integer of the prescribed parity, with a tie taken downward. That energy is either zero or greater than one half. Zero energy is exactly the case in which every phase is a positive integer of its prescribed parity.
\end{abstract}

\section{Words and energy}

Throughout, \(q\) and \(c\) are odd positive integers. For a bit \(e\in\{0,1\}\) define
\[
F_e(x)=
\begin{cases}
x/2 & \text{if }e=0,\\
(qx+c)/2 & \text{if }e=1.
\end{cases}
\]
A word is a finite sequence \(w=(e_0,\ldots,e_{m-1})\) with each \(e_i\in\{0,1\}\). Its length is \(m\), and \(r\) denotes the number of indices with \(e_i=1\). The word is strictly contracting when \(q^r<2^m\).

\begin{definition}
The affine data \((P(w),B(w))\) of a word are the positive integer, or nonnegative integer, defined as follows. The empty word has data \((1,0)\). If \(w\) is the bit \(e\) followed by a word \(w'\) with data \((P',B')\), then
\[
(P(w),B(w))=
\begin{cases}
(P',2B') & \text{if }e=0,\\
(qP',P'+2B') & \text{if }e=1.
\end{cases}
\]
\end{definition}

The first letter of \(w\) is applied first. Writing \(P=P(w)\) and \(B=B(w)\), the length-\(m\) composition is the affine map
\[
x\longmapsto \frac{Px+cB}{2^m}.
\]
Indeed, the empty composition is the identity, and appending a branch on the left multiplies the previous multiplier and offset by the corresponding rule. By induction on suffixes, \(P(w)=q^r\).

\begin{definition}
If \(q^r<2^m\), the fixed point of \(w\) and its phases are
\[
\xi_0=\frac{cB}{2^m-P},\qquad
\xi_{i+1}=F_{e_i}(\xi_i)\quad(0\le i<m).
\]
\end{definition}

\begin{lemma}
\label{lem:return}
If \(q^r<2^m\), then \(\xi_m=\xi_0\).
\end{lemma}

\begin{proof}
The length-\(m\) formula gives
\[
\xi_m=\frac{P\xi_0+cB}{2^m}.
\]
The denominator \(2^m-P\) is nonzero because \(P=q^r<2^m\). Substituting \(\xi_0=cB/(2^m-P)\) produces
\[
P\xi_0+cB
=P\xi_0+\xi_0(2^m-P)
=2^m\xi_0,
\]
so \(\xi_m=\xi_0\).
\end{proof}

\begin{definition}
\label{def:nearest}
Let \(x\in\Q\) and let \(e\in\{0,1\}\). Set \(k=\lfloor(x-e)/2\rfloor\) and \(\ell=e+2k\), and set
\[
\rho=
\begin{cases}
\ell & \text{if }x\le \ell+1,\\
\ell+2 & \text{if }x>\ell+1.
\end{cases}
\]
The selected integer \(n(e,x)\) is \(\rho\) when \(\rho>0\), and otherwise it is \(1\) when \(e=1\) and \(2\) when \(e=0\).
\end{definition}

Thus \(\ell\) is the unique integer of parity \(e\) satisfying \(\ell\le x<\ell+2\), and \(\rho\) is whichever of \(\ell\) and \(\ell+2\) lies closer to \(x\), with the smaller value chosen when \(x=\ell+1\).

\begin{lemma}
\label{lem:nearest}
The integer \(n=n(e,x)\) is positive and has parity \(e\). If \(m\) is any positive integer of parity \(e\), then \(\lvert x-n\rvert\le\lvert x-m\rvert\), and equality implies \(n\le m\). In particular, a positive integer of parity \(e\) is selected from itself.
\end{lemma}

\begin{proof}
The integer \(\ell\) has parity \(e\), and so do \(\ell+2\) and the positive replacements \(1\) and \(2\). Hence \(n\) has parity \(e\). Positivity holds because a nonpositive raw value is replaced by \(1\) or \(2\).

The two integers of parity \(e\) in an interval of length \(2\) are exactly \(\ell\) and \(\ell+2\). Every other positive integer \(m\) of parity \(e\) differs from \(\ell\) by \(2t\) for some nonzero integer \(t\). If \(t\ge 1\), then \(m\ge \ell+2\), so the distance from \(x\) to \(m\) is at least the distance from \(x\) to \(\ell+2\), which is at least the distance from \(x\) to \(\rho\). If \(t\le -1\), then \(m\le \ell-2\), so the distance from \(x\) to \(m\) is at least the distance from \(x\) to \(\ell\), again at least the distance to \(\rho\). Therefore \(\rho\) minimizes the distance among all integers of parity \(e\), and it is the smaller of the two minimizers when both are equally near.

If \(\rho>0\), then \(n=\rho\) and the same comparison applies to every positive parity integer. If \(\rho\le 0\), the sandwich \(\ell\le x<\ell+2\) forces \(x\le 1\). Let \(a\) be the least positive integer of parity \(e\): it is \(1\) for odd parity and \(2\) for even parity. The selector bound gives \(x\le a\). Every positive integer \(m\) of that parity satisfies \(a\le m\), so \(\lvert x-a\rvert=a-x\le m-x=\lvert x-m\rvert\), with equality only when \(m=a\). Equality of distances in either case forces the competitor to be at least \(n\), because a strictly smaller positive parity integer would lie strictly farther from \(x\).

Finally, if \(m\) itself is positive and has parity \(e\), the inequality with both sides equal to \(m\) forces \(\lvert m-n(e,m)\rvert=0\), so \(n(e,m)=m\).
\end{proof}

\begin{definition}
Assume \(q^r<2^m\). For \(0\le i<m\), let \(n_i=n(e_i,\xi_i)\). The carry energy of \(w\) is
\[
E(w)=\max_{0\le i<m}\lvert\xi_i-n_i\rvert,
\]
interpreted as \(0\) when \(m=0\). The hypothesis \(q^r<2^m\) excludes the empty word, since it would require \(1<1\).
\end{definition}

\section{The gap}

\begin{theorem}
\label{thm:carry-gap}
Let \(q\) and \(c\) be odd positive integers, and let \(w\) be a word of length \(m\) with \(r\) odd letters such that \(q^r<2^m\). Let \(E(w)\) be the carry energy of its fixed-point phases. Then
\[
E(w)=0\qquad\text{or}\qquad E(w)>\frac12.
\]
Furthermore, \(E(w)=0\) if and only if every phase \(\xi_i\) is a positive integer of parity \(e_i\).

For \(q=3\) and \(c=1\), the word \(10\) has energy \(0\), and the word \(100\) has energy \(8/5\).
\end{theorem}

\begin{proof}
Write \(a_i=q/2\) when \(e_i=1\) and \(a_i=1/2\) when \(e_i=0\). For \(0\le i<m\) set \(\eta_i=\xi_i-n_i\). Lemma~\ref{lem:return} gives \(\xi_m=\xi_0\). The return uses the first letter \(e_0\), so the selected integer at the return is \(n_0\) and the return error equals \(\eta_0\).

The branch identity is \(F_e(n+\eta)=F_e(n)+a\eta\), where \(a\) is the scale of \(e\). Since \(n_i\) has parity \(e_i\) and both \(q\) and \(c\) are odd, \(F_{e_i}(n_i)\) is an integer: an even integer is divisible by \(2\), and the sum of the two odd integers \(qn_i\) and \(c\) is even. Define the integer carry
\[
k_i=F_{e_i}(n_i)-n_{i+1},
\]
where \(n_m\) means \(n_0\). Then
\[
\eta_{i+1}=a_i\eta_i+k_i.
\]

Suppose \(E(w)\le 1/2\). Every phase error, including the return error \(\eta_0\), then has absolute value at most \(1/2\). If \(k_i=0\), the recurrence gives \(\lvert\eta_{i+1}\rvert=a_i\lvert\eta_i\rvert\). If \(k_i\neq 0\), then \(\lvert k_i\rvert\ge 1\). The triangle inequality yields
\[
\lvert k_i\rvert
\le\lvert\eta_{i+1}\rvert+\lvert a_i\eta_i\rvert,
\]
so
\[
\lvert a_i\eta_i\rvert
\ge\lvert k_i\rvert-\lvert\eta_{i+1}\rvert
\ge 1-\lvert\eta_{i+1}\rvert
\ge\lvert\eta_{i+1}\rvert,
\]
where the last step is \(\lvert\eta_{i+1}\rvert\le 1/2\). Thus \(\lvert\eta_{i+1}\rvert\le a_i\lvert\eta_i\rvert\) at every index.

The product of the scales \(a_0\cdots a_{m-1}\) equals \(q^r/2^m\). Iterating the comparison gives
\[
\lvert\eta_0\rvert
=\lvert\eta_m\rvert
\le\frac{q^r}{2^m}\lvert\eta_0\rvert.
\]
The coefficient is strictly less than \(1\), so \(\eta_0=0\). Each later error is then bounded by a scale product times zero, hence vanishes. Every absolute phase error is zero, and therefore \(E(w)=0\).

Consequently, the assumption \(E(w)\le 1/2\) forces \(E(w)=0\). If the energy is not at most \(1/2\), then \(E(w)>1/2\).

If \(E(w)=0\), then \(\xi_i=n_i\) for every \(i<m\). Lemma~\ref{lem:nearest} says that each \(n_i\) is a positive integer of parity \(e_i\). Conversely, if every phase is already a positive integer of parity \(e_i\), the same lemma says that the phase is selected from itself, so every distance vanishes and \(E(w)=0\).

For the word \(10\) with \(q=3\) and \(c=1\), the affine data are \((P,B)=(3,1)\) and the fixed point is \(1/(4-3)=1\). The phases are \(1\) and \(F_1(1)=2\). Both are positive integers of the prescribed parities, so the energy is \(0\).

For the word \(100\), the affine data are again \((3,1)\), now of length \(3\), and the fixed point is \(1/(8-3)=1/5\). The phases are
\[
\frac15,\qquad
F_1\Bigl(\frac15\Bigr)=\frac45,\qquad
F_0\Bigl(\frac45\Bigr)=\frac25.
\]
Definition~\ref{def:nearest} selects \(1\) for \(1/5\) and selects \(2\) for both \(4/5\) and \(2/5\). The absolute errors are \(4/5\), \(6/5\), and \(8/5\), so the energy is \(8/5\).
\end{proof}

\begin{remark}
Zero energy is the locus of positive integer phases of the prescribed parities; the theorem does not exclude it. The theorem does not assert that the classical Collatz map reaches \(1\), and it does not assert that energies arbitrarily close to \(1/2\) occur for the constant \(c=1\).
\end{remark}

\end{document}
